\section{Introduction}
\label{sec:intro}
Scaling up transformer-based language models has unlocked the in-context learning (ICL) paradigm, where the model can learn to perform a previously unseen task from a few demonstrated examples~\citep{brown2020language,chowdhery2023palm,ouyang2022training}.
However, improving ICL capabilities remains a challenging task~\citep{gendron2024largelanguagemodelsstrong,li2025mirageevaluatingexplaininginductive}.
Language models heavily rely on semantic priors from pretraining~\citep{wei2023larger}, which hurts their ability to perform inductive reasoning—a core feature of human intelligence~\citep{peirce1868questions}.
Inductive reasoning involves identifying general principles from specific examples and applying them to novel situations. In an ICL scenario, the language model must first generate a hypothesis that captures the input-output relationship in the demonstrations and then apply this hypothesis to a novel input to produce the desired output~\citep{wang2023hypothesis,qiu2023phenomenal}.
To achieve high accuracy, even very large models like GPT-4~\citep{openai2024gpt4technicalreport} often rely on advanced methods such as chain-of-thought (CoT)~\citep{li2025mirageevaluatingexplaininginductive}, hypothesis search~\citep{wang2023hypothesis}, and iterative hypothesis refinement~\citep{qiu2023phenomenal}.
Although the recent reasoning model DeepSeek-R1~\citep{guo2025deepseek} shows impressive reasoning ability via test-time scaling~\citep{openai2024openaio1card}, few-shot prompting consistently degrades its performance, indicating that ICL remains a challenge.

Although instruction-tuning can be used to improve the ICL capabilities of a language model, previous works in this area mostly rely on imitation learning~\citep{huang2022incontextlearningdistillationtransferring,alpaca,vicuna2023}. Imitation learning teaches a student model to directly mimic the outputs of a teacher model but fails to teach the underlying rules connecting inputs and outputs in complex reasoning tasks~\citep{mukherjee2023orcaprogressivelearningcomplex}. 
In this work, we aim to improve the ICL capabilities of language models on inductive reasoning tasks via a novel reasoning distillation technique. We propose \proposed, a reasoning distillation method specifically designed to enhance the inductive reasoning abilities of language models. 
The main challenge in reasoning distillation is that the underlying rules are often not available as ground truth. We address this issue by using a teacher model to generate a set of candidate hypotheses to approximate the underlying rules. We also demonstrate that noisy fitness estimation—using natural language rule-following with the teacher model—can be effectively leveraged to evaluate the quality of rules, thereby improving model tuning.
\proposed consists of three key steps: (1) a data augmentation step, where we generate datasets for both rule generation and rule-following components of inductive reasoning using a teacher model; (2) a supervised fine-tuning (SFT) step, where the student model learns to mimic the rule generation and rule-following behavior of the teacher; and (3) an alignment step, which prepares the student model for inference-time optimizations for advanced inductive reasoning methods such as hypothesis search.

We find that \proposed helps small models like LLaMA-3, Mistral, and Qwen outperform proprietary models on inductive reasoning tasks, and even surpass the teacher model, GPT-4o~\citep{openai2024gpt4ocard}, in some cases.
We investigate the key steps and factors affecting the distillation process. Specifically, we show that the quality of a generated hypothesis can be evaluated by checking the number of demonstrated ICL training examples that satisfy that hypothesis. This measurement allows us to create a preference alignment corpus for rule generation and to filter the SFT corpus, both of which positively impact performance.
Finally, we examine the effects of different alignment techniques and observe that Odds Ratio Preference Optimization (ORPO)~\citep{hong2024orpomonolithicpreferenceoptimization} achieves the best overall performance compared to Direct Preference Optimization (DPO)~\citep{rafailov2024directpreferenceoptimizationlanguage} and Kahneman-Tversky Optimization (KTO)~\citep{ethayarajh2024ktomodelalignmentprospect} in our setting. We find that the preference alignment stage can significantly improve the inference-time efficiency of inductive reasoning.

Our contributions are as follows.
\begin{itemize}
    \item We propose \proposed, a reasoning distillation method specifically designed to improve the inductive reasoning ability of language models. \proposed demonstrates enhanced performance on four diverse tasks: 1D-ARC~\citep{xu2024llmsabstractionreasoningcorpus}, List Function~\citep{rule2020child}, ACRE~\citep{zhang2021acreabstractcausalreasoning}, and MiniSCAN~\citep{lake2019humanfewshotlearningcompositional}. Notably, the proposed approach outperforms even the teacher model, GPT-4o, on three out of four tasks within an equivalent hypothesis search space. \proposed, based on the LLaMA-3 backbone, achieves relative improvements of 23.2\%, 2.8\%, and 66.6\% over GPT-4o on 1D-ARC, ACRE, and MiniSCAN, respectively.
    
    \item We demonstrate the impact of different combinations of model tuning and alignment techniques on overall performance. We also perform inference-time optimization by aligning the models to generate higher-quality hypotheses within a smaller candidate search space compared to the teacher. \proposed, based on the LLaMA-3 backbone, performs equal to or better than the teacher model while being 87\%, 53\%, and 25\% more token-efficient on MiniSCAN, 1D-ARC, and ACRE, respectively.
\end{itemize}

